% Fichier complété par tools/complete_missing_corrections.py.
% Source : COR/COR-02-P/65_Eclipse/65_Eclipse.tex
% Statut : rédigé (3 question(s)).
\begin{corrigebox}[Corrigé rédigé]
\CorrigeQuestion{1}
On calcule l'erreur finale par la valeur finale, le temps de réponse sur
            la réponse indicielle et les marges sur le Bode/Nyquist. Le système non corrigé
            présente le compromis habituel : augmenter $K_P$ réduit l'erreur et accélère
            la réponse mais diminue les marges de stabilité. Les critères sont comparés un
            par un aux seuils du cahier des charges.
\CorrigeQuestion{2}
Non en général : une unique action proportionnelle ne permet pas de régler
            indépendamment précision, rapidité et stabilité. La valeur imposée par la
            précision conduit à une marge insuffisante, tandis que la valeur assurant la
            stabilité laisse une erreur trop grande; une action intégrale/avance de phase
            est nécessaire.
\CorrigeQuestion{3}
Une perturbation en échelon est rejetée sans erreur permanente seulement
            si la fonction de transfert de la boucle vue depuis la perturbation possède un
            gain infini en basse fréquence. Une correction purement proportionnelle ne
            l'assure pas : le système n'est donc pas robuste au sens de l'erreur nulle.
\end{corrigebox}
